\section{Métodos}
\subsection{Diseño y procedencia de las entradas}
Se realizó un cálculo matemático de reconstrucción retrospectiva. Se utilizaron las masas de reposo y demás propiedades del inventario HRG ya empleado, y se añadió la abundancia cosmológica actual como condición para inferir las poblaciones tempranas. Se calcularon por separado las densidades de bariones y antibariones, su entropía y el exceso neto. Las comparaciones de entropía a potenciales químicos nulos se conservaron como controles del inventario. Se evaluaron las temperaturas
\[
T=150,\ 140,\ 120,\ 100,\ 80,\ 60,\ 50,\ 40\ {\rm MeV}
\]
La reconstrucción resolvió $\mu_B$ en cada temperatura, manteniendo $\mu_Q=\mu_S=0$; el control simétrico fijó también $\mu_B=0$. Aquí $\mu_B$ regula el desequilibrio entre bariones y antibariones, y $\mu_Q,\mu_S$ son los potenciales asociados a carga eléctrica y extrañeza. Se adoptaron unidades naturales, $\hbar=c=k_B=1$, y se convirtieron las temperaturas y masas a GeV antes de realizar las sumas. Todas las especies de partícula compartieron la misma temperatura. Se denominó especie de partícula a cada tipo definido por una entrada del inventario; $g_i$ contó sus estados de espín. El punto inicial $T_i=150$ MeV fue el extremo superior de la malla, sin asignarle una edad cósmica ni identificarlo con el inicio cinético de aniquilación.

Se utilizaron exclusivamente los archivos \texttt{list.dat} de los directorios PDG2014, PDG2020 y PDG2025 del repositorio Thermal-FIST \citep{TFLists2026}. Las entradas quedaron fijadas en la revisión \texttt{4086191e9522}; el suplemento registra el identificador completo, enlaces exactos y huellas SHA-256. Las listas seleccionadas contienen hadrones del sector de quarks $u,d,s$. No se añadieron las listas separadas de hadrones con quark \emph{charm}, de núcleos o de otros sectores; no representan la totalidad de todas las partículas del PDG.

La Tabla~\ref{tab:inventory} identifica las dimensiones efectivas de cada entrada. Se descartaron líneas comentadas y contribuciones con degeneración cero. Se conservaron las masas y degeneraciones de los archivos, incluidas las asignaciones de espín adoptadas por sus responsables.

\begin{table}[t]
\centering\small
\caption{Inventario de las listas utilizadas. Las dos últimas columnas incluyen antipartículas y sólo entradas con degeneración positiva.}
\label{tab:inventory}
\setlength{\tabcolsep}{4pt}
\begin{tabular}{lrrr}\toprule
Lista & Filas & Especies activas & $B+\bar B$\\\midrule
PDG2014 & 219 & 387 & 260\\
PDG2020 & 242 & 433 & 298\\
PDG2025 & 248 & 443 & 292\\
\bottomrule\end{tabular}
\end{table}

Los valores reducidos se transcribieron de \citet{Garcia2026} y se almacenaron como una entrada independiente, con cuatro cifras significativas. No se reconstruyó una lista reducida ad hoc para forzar esa tabla. Las cifras de las columnas PDG que se presentan más adelante son resultados de este trabajo calculados a partir de datos publicados; no fueron extraídas como esos mismos porcentajes de una publicación externa.

\subsection{Observable y balance termodinámico}
Se definió la fracción instantánea
\begin{equation}
\fB(T)=\frac{\sB(T)}{\stot(T)},
\label{eq:f}
\end{equation}
donde $\sB$ incluye todos los bariones y antibariones activos de la lista elegida. A potenciales químicos nulos, este sector es simétrico. La entropía hadrónica se escribió como
\begin{equation}
\shad=\sB+s_{\rm mes},
\end{equation}
y la del plasma considerado como
\begin{equation}
\stot=\shad+s_\gamma+s_{e^\pm}+s_{\mu^\pm}+s_{\nu\bar\nu}.
\label{eq:den}
\end{equation}
Para comprobar el efecto del denominador se calculó también
\begin{equation}
f_{\rm had}(T)=\frac{\sB(T)}{\shad(T)}.
\label{eq:fhad}
\end{equation}
Ninguna de estas fracciones es una abundancia neta, una tasa de aniquilación ni una fracción de los fotones actuales con ascendencia identificable.

\subsection{Entropía por especie de partícula y conteo de antipartículas}
Para cada especie de partícula masiva de un gas ideal de equilibrio, a $\mu_i=0$, se evaluó
\begin{align}
s_i(T)&=\frac{g_i m_i^2T}{2\pi^2}
\sum_{k=1}^{K_{\max}}c_{i,k}\,A_{i,k},\label{eq:entropy}\\
A_{i,k}&=\frac{4}{k^2}K_2(kx_i)+\frac{x_i}{k}K_1(kx_i),\\
x_i&=\frac{m_i}{T}.\nonumber
\end{align}
$K_1$ y $K_2$ son funciones de Bessel modificadas de segunda clase; $g_i$ es la degeneración de espín. Se empleó $c_{i,k}=1$ para bosones y $c_{i,k}=(-1)^{k+1}$ para fermiones. Estos coeficientes corresponden a las distribuciones de Bose-Einstein para bosones y Fermi-Dirac para fermiones, que describen cómo se ocupan los estados de energía a la temperatura elegida. Retener sólo $k=1$ recupera la aproximación de Maxwell-Boltzmann. Se trata de distribuciones físicas de población, distintas de las pruebas estadísticas sobre mediciones.

Cada archivo contiene un miembro de los pares partícula-antipartícula. Se multiplicó su contribución por dos cuando al menos uno de sus números cuánticos $B,Q,S,C$ era no nulo, y por uno para estados autoconjugados. Los estados de carga o isoespín ya figuran como entradas individuales, por lo que no se añadió otra degeneración de isoespín. La implementación verificó la ausencia de identificadores PDG negativos y de duplicados antes de aplicar esta regla.

Se usaron masas nominales y anchos de resonancia nulos. Las tablas de desintegración no se propagaron ni se añadieron de nuevo al denominador: el objeto calculado fue la entropía del inventario de equilibrio, no los rendimientos finales después de las desintegraciones. Tampoco se incluyeron correcciones de volumen excluido, interacciones atractivas o repulsivas ni ajustes a datos de red.

\subsection{Reconstrucción del número de partículas}
Además de $s_i$, se calculó la densidad numérica de cada especie de partícula mediante las distribuciones térmicas del gas ideal \citep{Vovchenko2019}:
\begin{equation}
n_i^{\rm eq}(T)=\frac{g_i m_i^2T}{2\pi^2}
\sum_{k=1}^{96}\frac{c_{i,k}}{k}K_2(km_i/T).
\label{eq:number}
\end{equation}
Aquí $g_i$ correspondió a una especie de partícula; su antipartícula se registró por separado con igual densidad a $\mu_i=0$. Se sumaron $n_B$ sobre $B_i=+1$ y $n_{\bar B}$ sobre $B_i=-1$. Se verificó que todas las entradas bariónicas activas tenían $|B_i|=1$. Se definieron
\begin{align}
n_\Sigma&=n_B+n_{\bar B}, & \Delta n_B&=n_B-n_{\bar B},\\
Y_\Sigma&=n_\Sigma/\stot, & \YB&=\Delta n_B/\stot.
\label{eq:yields}
\end{align}
El control con $\mu_B=0$ dio $\Delta n_B=0$; la reconstrucción con exceso finito se describe más adelante. Como contabilidad de pares potenciales se usó $n_{\rm par}=\min(n_B,n_{\bar B})$; no representa pares ligados ni eventos de reacción contabilizados.

Se convirtieron las densidades naturales a ${\rm cm^{-3}}$ dividiendo entre $(\hbar c)^3$, con $\hbar c=1.973269804\times10^{-14}$ GeV cm \citep{Robinson2025}. El número medio en un volumen especificado es $N_i=n_iV$. Se exportó el inventario de densidades por especie de partícula y temperatura. La reconstrucción consistió en inferir estas poblaciones de equilibrio mediante sus propiedades y la condición cosmológica; una regeneración física requeriría integrar tasas de producción inversa y aniquilación.

\subsection{Sector ligero común}
Los fotones se trataron como bosones sin masa con $g=2$; electrones y positrones, como fermiones sin masa con $g=4$ en conjunto; y los tres sabores de neutrinos y antineutrinos, como fermiones sin masa con seis estados de helicidad en conjunto. Para cada componente relativista,
\begin{equation}
s_{\rm rel}=\frac{2\pi^2}{45}gT^3
\begin{cases}
1,&\text{bosones},\\
7/8,&\text{fermiones}.
\end{cases}
\end{equation}
Los muones y antimuones se evaluaron con la ecuación~\eqref{eq:entropy}, $g=4$ en conjunto y $m_\mu=0.1056583755$ GeV \citep{PDG2025}. Este mismo sector ligero se utilizó para las tres listas. La aproximación de electrones sin masa se comprobó por integración con su masa física y se cuantificó por separado.

\subsection{Normalización al residuo y a las abundancias cosmológicas}
Se adoptaron como entradas externas $\eta=(6.12\pm0.04)\times10^{-10}$ y $\YB^{\rm ref}=\eta/7.04$ \citep{Fields2025}. El factor 7.04 correspondió a la aproximación estándar de entropía por fotón después de la aniquilación $e^+e^-$. Para comparar épocas con distinta dilución se emplearon los rendimientos $Y$, en lugar de comparar directamente densidades físicas. Se calculó
\begin{equation}
\mathcal Q(T)=\frac{Y_\Sigma^{\rm eq}(T)}{\YB^{\rm ref}}.
\label{eq:populationratio}
\end{equation}
$\mathcal Q$ es el número térmico total por unidad de exceso neto de referencia; $\mathcal Q/2$ es su contabilidad simétrica de pares. No es una fracción de supervivencia medida.

Para conectar esa normalización con el presente se usó $T_{\gamma0}=2.7255\pm0.0006$ K, valor redondeado de la determinación de \citet{Fixsen2009}. Se obtuvo
\begin{align}
n_{\gamma0}&=\frac{2\zeta(3)}{\pi^2}
\left(\frac{k_BT_{\gamma0}}{\hbar c}\right)^3,\\
s_0/k_B&\simeq7.04\,n_{\gamma0}, & n_{b0}&=\eta\,n_{\gamma0}.
\label{eq:today}
\end{align}
$n_{b0}$ contó unidades de número bariónico (nucleones), incluidas las contenidas en núcleos; no contó átomos. En una comparación adiabática, $s(T)V(T)=s_0V_0$, de modo que $N_\Sigma^{\rm eq}=Y_\Sigma^{\rm eq}s_0V_0$ y $N_{\Delta B}^{\rm ref}=\YB^{\rm ref}s_0V_0$. Esto estableció una normalización comóvil común, condicionada a conservación de entropía y número bariónico neto.

Con la fracción bariónica primordial en $^4$He, $Y_p=0.245\pm0.003$, y $(D/H)_{\rm P}=(2.508\pm0.029)\times10^{-5}$ \citep{Fields2025}, se calculó una composición primordial equivalente a la densidad media actual:
\begin{align}
n_{{\rm H},0}^{\rm prim}&\simeq(1-Y_p)n_{b0},\\
n_{{}^4{\rm He},0}^{\rm prim}&\simeq Y_p n_{b0}/4,\\
n_{{\rm D},0}^{\rm prim}&\simeq(D/H)_{\rm P}n_{{\rm H},0}^{\rm prim}.
\label{eq:hydrogen}
\end{align}
Se omitieron las pequeñas correcciones de enlace nuclear y de abundancia de núcleos traza. Estas conversiones no incorporaron procesamiento estelar ni fracción de ionización; por ello no describen una medición de hidrógeno neutro actual. Se propagaron a primer orden las incertidumbres de $\eta,T_{\gamma0},Y_p,D/H$, tratándolas como independientes. Para cualquier salida $F$ se usó $\sigma_F^2\simeq\sum_j(\partial F/\partial x_j)^2\sigma_{x_j}^2$; las salidas comparten entradas y sus errores están correlacionados. No se atribuyeron esas incertidumbres al modelo HRG.

\subsection{Inferencia inversa desde la abundancia actual}
La condición central fue conservar el rendimiento neto bajo expansión adiabática y sin violación del número bariónico:
\begin{equation}
\frac{\Delta n_B(T,\alpha)}{\stot(T,\alpha)}
=\frac{\eta}{7.04},\qquad \alpha=\mu_B/T.
\label{eq:inverse}
\end{equation}
Se utilizó el valor central de $\eta$ para resolver esta condición en cada lista y temperatura. Su incertidumbre observacional no se propagó a las poblaciones inversas; esas cifras son estimaciones condicionadas a las entradas centrales.

Para las entradas con $B_i=+1$, se definieron $a_{i,k}=g_i m_i^2T c_{i,k}K_2(kx_i)/(2\pi^2 k)$ y $b_{i,k}=g_i m_i^2T c_{i,k}A_{i,k}/(2\pi^2)$. Las antipartículas tuvieron potencial opuesto. Se evaluaron
\begin{align}
n_\Sigma(T,\alpha)&=2\sum_{i,k}a_{i,k}\cosh(k\alpha),\\
\Delta n_B(T,\alpha)&=2\sum_{i,k}a_{i,k}\sinh(k\alpha),\\
\sB(T,\alpha)&=2\sum_{i,k}\bigl[b_{i,k}\cosh(k\alpha)\nonumber\\
&\hspace{18mm}-\alpha a_{i,k}\sinh(k\alpha)\bigr].
\label{eq:inversekernels}
\end{align}
Esta forma evitó restar dos densidades casi iguales para obtener el pequeño exceso. Se añadió la entropía mesónica y ligera y se resolvió la ecuación~\eqref{eq:inverse} mediante el método de Brent, en $0\leq\alpha\leq1$, con tolerancias absoluta $10^{-18}$ y relativa $10^{-13}$. Se recuperaron $n_B=(n_\Sigma+\Delta n_B)/2$ y $n_{\bar B}=(n_\Sigma-\Delta n_B)/2$.

Las masas, el espín y las cargas provinieron del catálogo; $\eta$ y la temperatura actual del CMB fueron entradas observacionales. La temperatura temprana fue una condición elegida, y el HRG ideal aportó la relación térmica necesaria para inferir el total. No se ajustaron masas ni se analizaron nuevas firmas espectrales. Mantener $\mu_Q=\mu_S=0$ no impuso simultáneamente neutralidad eléctrica y extrañeza neta nula; este alcance se conservó al interpretar la reconstrucción.

\subsection{Temperatura inicial y referencia física}
Se conservó $T_i=150$ MeV como inicio de la malla principal para mantener la comparación con nuestros cálculos previos. Para fundamentar su escala, se utilizó la estimación de transición de QCD de \citet{Bazavov2019}, $T_c=156.5\pm1.5$ MeV, obtenida mediante cálculos en red a potenciales químicos nulos. Esa referencia caracteriza la aparición de hadrones, sin fijar un inicio universal de aniquilación. Se repitió la solución inversa a 150, 155, 156.5 y 158 MeV con las tres listas y con 96 y 192 términos. Se convirtió a kelvin mediante $T_{\rm K}=10^6T_{\rm MeV}/k_B$, con $k_B=8.617333262\times10^{-5}$ eV K$^{-1}$ \citep{Robinson2025}. La evaluación cerca de $T_c$ se trató como sensibilidad del HRG ideal; la incertidumbre publicada de $T_c$ no se interpretó como error físico completo de nuestras densidades.

\subsection{Prueba de unicidad y contraste con la literatura}
Se examinó la información de la abundancia neta mediante el mapa lineal de rendimientos
\begin{equation}
Y_{\Delta B}=\begin{pmatrix}1&-1\end{pmatrix}
\begin{pmatrix}Y_B\\Y_{\bar B}\end{pmatrix}.
\label{eq:rank}
\end{equation}
Se calculó su rango y su dirección nula. Añadir una cantidad $q$ a ambos componentes conserva la diferencia. Se construyeron cuatro ejemplos algebraicos con totales iguales a 0.1, 1, 10 y 100 veces el total reconstruido a 150 MeV, manteniendo el mismo exceso. Esos ejemplos examinaron la unicidad bajo conservación; no demostraron historias cinéticas ni estados de equilibrio físicamente realizables a una misma temperatura.

Se evaluó además el límite de nucleones de Maxwell-Boltzmann que sustenta la ecuación (19) de \citet{Satarov2013}, usando las masas distintas de protón y neutrón de nuestra lista:
\begin{equation}
n_N^{\rm MB}(T)=\sum_{i=p,n}
\frac{g_i m_i^2T}{2\pi^2}K_2(m_i/T).
\label{eq:nucleonMB}
\end{equation}
Se comparó con la suma de protones y neutrones obtenida con Fermi-Dirac y con todos los bariones del HRG, a $\mu_B=0$. Se calcularon la participación nucleónica y la diferencia relativa entre las dos distribuciones. Este contraste fue nuestra evaluación de una aproximación publicada; no reprodujo el cálculo cinético con interacciones de esos autores ni utilizó sus poblaciones como mediciones.

\subsection{Procedimiento ejecutado paso a paso}
La reconstrucción y sus controles se realizaron en el siguiente orden:
\begin{enumerate}
\setlength{\itemsep}{2pt}
\item Se leyeron las tres listas fijadas, se extrajeron masas de reposo, espín, distribución térmica y cargas, y se verificó el inventario.
\item Se calcularon $s_i$ y $n_i$ a $\mu_B=0$ con 96 términos y se sumaron bariones, antibariones, mesones y el sector ligero común.
\item Se convirtió la abundancia publicada $\eta$ a $Y_{\Delta B}^{\rm ref}$ y se resolvió la condición inversa en cada temperatura para obtener $\alpha$, $n_B$ y $n_{\bar B}$.
\item Se verificaron el rango del mapa neto y la existencia de poblaciones algebraicas distintas con el mismo exceso.
\item Se compararon el límite nucleónico publicado, el HRG completo y la referencia reducida de entropía, sin ajustar parámetros para mejorar la concordancia.
\item Se repitieron las sumas con 192 términos y se contrastaron especies representativas y excesos a potencial finito con integrales independientes.
\item Se calcularon las conversiones cosmológicas y sus incertidumbres. Se exportaron resultados y procedencia, y se generaron las tablas y figuras desde esos archivos.
\end{enumerate}

\subsection{Análisis descriptivos y procedimientos estadísticos}
Con $f_{\rm previo}$ como referencia reducida y $f_{25}$ como el cálculo PDG2025, se definieron
\begin{align}
R(T)&=\frac{f_{\rm previo}(T)}{f_{25}(T)},\\
d(T)&=1-R(T),\\
\Delta_{\rm pp}(T)&=100\,[f_{25}(T)-f_{\rm previo}(T)].
\end{align}
$100R$ y $100d$ se expresan en porcentaje; $\Delta_{\rm pp}$ es una diferencia en puntos porcentuales de la fracción física. El término \emph{déficit relativo} describe esta diferencia de modelos, sin identificar una cantidad de partículas o de entropía materialmente perdida.

Para las ocho temperaturas se calcularon la correlación de Pearson entre $\log_{10}f_{\rm previo}$ y $\log_{10}f_{25}$, la correlación de rangos de Spearman y el error cuadrático medio logarítmico:
\begin{equation}
{\rm RMSE}_{\log}=
\sqrt{\frac{1}{8}\sum_{j=1}^{8}
\left[\log_{10}\frac{f_{\rm previo}(T_j)}{f_{25}(T_j)}\right]^2}.
\end{equation}
La amplitud entre catálogos se midió como $100(\max f-\min f)$ en cada temperatura. Es una sensibilidad a esas tres entradas correlacionadas, no un intervalo estadístico de incertidumbre. Pearson midió similitud lineal en escala logarítmica; Spearman, orden monotónico; RMSE cuantificó discrepancia de amplitud. Se emplearon como estadísticos descriptivos. Las ocho temperaturas fueron puntos calculados con entradas fijas, sin réplicas experimentales; no se realizaron pruebas de hipótesis con valores $p$, pruebas $t$ o ANOVA ni ajustes cosmológicos de verosimilitud.

\subsection{Verificación numérica}
Las sumas de 96 y 192 términos se compararon para las tres listas y todas las temperaturas, exigiendo una diferencia relativa inferior a $10^{-10}$ en $\fB$ y en las densidades $n_i$. Para $\pi^0,\pi^\pm,p,n$ y $\mu^\pm$ se integró de forma independiente
\begin{equation}
s_i=\frac{g_i}{2\pi^2T}\int_0^\infty
p^2\left(E_i+\frac{p^2}{3E_i}\right)
\frac{dp}{e^{E_i/T}+\epsilon_i},
\end{equation}
con $E_i=\sqrt{p^2+m_i^2}$ y $\epsilon_i=+1$ para fermiones, $-1$ para bosones. Se comprobaron 150 y 40 MeV con un criterio relativo de $10^{-8}$. Para la densidad numérica se contrastó de igual modo
\begin{equation}
n_i=\frac{g_i}{2\pi^2}\int_0^\infty
\frac{p^2\,dp}{e^{E_i/T}+\epsilon_i}.
\label{eq:numberintegral}
\end{equation}
Se verificaron las identidades de suma y conjugación, positividad, cotas de las fracciones y monotonicidad en la malla. Las integraciones se realizaron con cuadratura adaptativa, tolerancias absoluta $10^{-15}$ y relativa $5\times10^{-12}$ en las integrales adimensionales; se contrastaron los resultados con un criterio relativo de $10^{-8}$.

Para la reconstrucción inversa se exigió un residuo relativo inferior a $10^{-10}$ en la ecuación~\eqref{eq:inverse} y se repitió la solución con 192 términos. Se comprobaron por cuadratura independiente los excesos de protones y neutrones a 150 y 40 MeV. El integrando adimensional fue $q^2\sinh\alpha/[\cosh(\sqrt{q^2+x_i^2})+\cosh\alpha]$, multiplicado por $g_iT^3/(2\pi^2)$, con $q=p/T$. Se mantuvieron las tolerancias anteriores y el criterio relativo de $10^{-8}$. También se verificaron $n_B>n_{\bar B}>0$ y la recuperación del cálculo simétrico al fijar $\alpha=0$.
